% !TeX root = ../main.tex
\subsection{Model and inference}

For guide $g$ in library $i$, BARCS models the count conditional on the
unfiltered mapped-guide total:
\begin{equation}
K_{gi}\mid N_i\sim
\operatorname{BetaBinomial}(N_i,\mu_{gi},\rho_g),
\qquad
\logit(\mu_{gi})=\mathbf{x}_i^\top\boldsymbol{\beta}_g .
\label{eq:barcs-model}
\end{equation}
The coefficient model in Equation~\eqref{eq:barcs-model} allows the design
vector to contain time, dose, donor, batch, treatment, and
interaction terms.  Coefficients are estimated by beta-binomial weighted
logistic regression.  A prespecified coefficient or contrast is divided by
its model-based standard error and compared with a Student $t$ distribution
using residual sample degrees of freedom.  All reported tests are two-sided,
and guide- or gene-level probabilities are adjusted by Benjamini--Hochberg.
The guide-specific overdispersion $\rho_g$ is estimated by solving the Pearson
equation at the fitted mean, truncated at zero when the counts are not
overdispersed, and the coefficient and dispersion updates are alternated to
convergence.  BARCS-ST reports these unmoderated statistics.  BARCS-MOD
additionally shrinks each guide's fitted variance toward a trend in log mean
abundance by an empirical-Bayes step that changes standard errors and residual
degrees of freedom but not the coefficient estimates; the real-data benchmarks
use unmoderated BARCS.  The full estimating equations are given in the
Supplement.

\subsection{Gene-level aggregation and control calibration}

Guide coefficients are combined into a gene-level statistic by a prespecified
directional Stouffer summary: each guide's two-sided probability is converted
to a normal deviate carrying the sign of its coefficient, the deviates for a
gene are combined by the Stouffer rule, and the result is returned to a
two-sided gene probability, which is then adjusted by Benjamini--Hochberg.  The
reported gene effect is the median guide coefficient.

Control calibration rescales a reported statistic without altering the
coefficient estimates or their ranking.  Given a set of negative-control
features, \texttt{bb\_calibrate\_controls()} maps the 95th percentile of the
absolute signed-normal control statistic onto the two-sided standard-normal
0.05 cutoff and divides by the resulting scale, with scales below one truncated
at one so that the adjustment can only be conservative.  Two properties matter
for interpretation.  First, the scale must be estimated at the same aggregation
level as the statistic it corrects, because a scale learned from individual
control guides does not transfer to multi-guide gene statistics; non-targeting
guides are therefore grouped into pseudo-genes matching the target-gene
guide-count distribution before the rule is applied to a gene-level summary.
Second, the controls used to estimate the scale cannot also be used to evaluate
it.  Wherever a control-based error rate is reported, the scale is obtained by
deterministic five-fold cross-fitting and each fold is evaluated with the scale
estimated from the other four.

\paragraph{Validity of gene-level probabilities.}
The directional Stouffer summary treats the guides of a gene as independent
sources of evidence.  Guides targeting the same gene are generally positively
correlated---through a shared target effect, a shared cutting response,
copy-number-associated toxicity, clonal structure, or unmodeled library
composition shifts---and under positive correlation the combined statistic is
anti-conservative; the null-calibration grid in the Supplement quantifies the
resulting gene-level error.  Gene-level probabilities are therefore reported as
empirically calibrated ranking statistics rather than as tests carrying a
nominal false-discovery guarantee.  Guide-level coefficients, their standard
errors, and their ordering are unaffected by this limitation, and gene-level
inference is recovered when the gene-level null is established by permutation
that preserves the within-gene guide structure and the experimental blocking.

\subsection{Denominator choice defines the estimand}

The two denominators available in BARCS answer different questions and are not
interchangeable.  Under the full-library total, $\beta$ is the change in a
guide's share of the sequenced library.  This quantity is always well defined,
but it coincides with the change in absolute guide abundance only when the
overall library composition is stable across the design; when a large fraction
of guides deplete, every surviving guide's share rises and true-null guides
acquire a positive shift (Figure~\ref{fig:simcrispr}).  Under a control-set
denominator, $\beta$ is the change relative to that control set, and it
identifies the absolute effect under the assumption that the control set is
invariant across design levels.  That assumption should be checked rather than
assumed, by comparing the control set's share of the library across design
levels.  Non-targeting controls are the efficient choice when the quantity of
interest is a targeted effect; safe-harbor controls additionally absorb the
shared cutting response and are preferred in Cas9 knockout screens with
widespread depletion.

\subsection{Longitudinal Cas13 analysis}

For HAP1, HEK293FT, MDA-MB-231, and THP1, days 0, 7, and 14 were fitted
jointly with continuous time and a replicate indicator.  K562 was excluded
because one day-0 replicate was absent.  BARCS, MAGeCK-MLE, edgeR-QL, DESeq2,
and limma--voom received the identical rounded matrix and tested the same time
slope.  Non-targeting guides were grouped into pseudo-genes matching the
target-gene guide-count distribution, and the same post-aggregation control
scaling rule was applied to all five methods.  A BARCS endpoint fit using only
days 0 and 14 tested the value of the intermediate time point.  Known
essential genes were positives; targeted lncRNAs with TPM equal to zero in
both available expression assays were proxy nulls.  That definition is
imperfect in a direction that favours conservative methods, so the proxy-null
rates are diagnostics rather than validated type-I errors.  Because the deposited values were already normalized and
ComBat-corrected, this analysis was treated as a processed-count sensitivity
study.

\subsection{Ordered-bin IL2RA analysis}

The IL2RA screen comprised four ordered FACS bins from each of three donors.
BARCS fitted one ordered-bin score with donor indicators, and MAGeCK-MLE
received the same design.  Of 593 non-targeting guides, four folds estimated
the control scale and the omitted fold evaluated it; this was repeated across
five deterministic folds.  A seeded permutation of fold assignments checked
that identifier ordering did not determine the result.  Method discoveries
were compared with 26 regulators
validated by individual knockout and flow cytometry.  Waterbear and MAUDE
results were taken from their published analyses of the same experiment.

\subsection{Simulation benchmarks}

CRISPulator generated three 10,000-gene FACS screens with five guides per gene
and four independent replicates.  BARCS-ST and BARCS-MOD differed only by
guide-dispersion moderation.  simCRISPR generated three factorial
knockout-by-treatment screens with 2,000 sgRNAs and three replicates per arm.
Its known interaction effects allowed direct comparison of full-library,
non-targeting, and safe-harbor denominators.  Control guides used for
normalization and calibration were held out from scoring.

\subsection{External and supplementary analyses}

The external false-discovery audit re-evaluated the deposited Avana
CB\textsuperscript{2} probabilities using full-library totals and ran BARCS on
the deposited PSN1 and DeWeirdt normalized tables.  Supplementary analyses
include the serial-harvest HT-29 screen, a four-coefficient GSE70038 screen
with method-specific pathway enrichment, and the Sanson A375 endpoint and
copy-number audit.  The Supplement reports their complete preprocessing,
design, and scoring definitions in the same order as the corresponding
results.
